\documentclass{uerj}
\usepackage[utf8]{inputenc}
\usepackage[T1]{fontenc}
\usepackage{geometry}
\usepackage{graphicx}
\usepackage{float}
\usepackage{hyperref}
\usepackage[english,indonesian]{babel}

\title{Optimasi Pemrosesan Template LaTeX Menggunakan Pendekatan Berbasis AI}
\author{Academic Support Agent}
\date{\today}

\begin{document}
\maketitle

\begin{abstract}

\end{abstract}

\chapter{PENDAHULUAN}


\chapter{TINJAUAN PUSTAKA}


% Original content (commented out):
% % Modelo de Dissertação em Latex para o PPG em Engenharia Mecânica da UERJ
% Este modelo foi adaptado da versão disponibilizada no site da Engenharia Elétrica da UERJ
% http://www.pel.uerj.br/publico/Modelo_LaTeX_Dissertacao_UERJ.rar
% http://www.pel.uerj.br/defesas/
%
% Utilizei o WinEdt 6.0 com Miktex 2.9
%
% Para gerar o PDF usei a opção PDFTeXfy com o documento [masterthesis.tex] aberto e em foco.
%
% Não consegui usar as figuras em EPS como o modelo original. Usei PNG e JPG sem problemas.
%
% Felipe M. - 20/06/2012
%

% Modificações e adaptações ao modelo da UFPE feitas por João Cordeiro.
% Versão original disponibilizada como template do latex.
% Este modelo foi criado para servir de exemplo de formato, não devendo seu código ser
% utilizado como referência de formatação.
% PLAIN: O código é uma adaptação de uma adaptação. Funciona, mas está cheio de gambiarras.
%
% João Cordeiro - 05/03/2019
%

% Atualização do template com outros pacotes e resolução de Erros no código.
% Palloma Marciano - 26/05/2021. (Não esquecer de fazer o upload da BiBTeX

%\documentclass[a4paper,12pt,oneside,openany,table,xcdraw]{uerj}
%\documentclass[a4paper]{report}
\documentclass[12pt,oneside,a4paper,brazil,]{abntex2}
%\documentclass[a4paper,12pt,oneside,openany,table,xcdraw]{uerj}

%\usepackage[portugues]{babel}
%\usepackage[latin1]{inputenc}
\usepackage[utf8]{inputenc}
\usepackage{enumerate}
\usepackage{cite}
\usepackage{epsf,epsfig,psfig}
\usepackage{pagina}
\usepackage{indentfirst}
\usepackage{theorem}
\usepackage{fancyhdr}
\usepackage{setspace}
\usepackage{boxedminipage}
\usepackage{float}
\usepackage{makeidx}
\usepackage{amsmath}
\usepackage{hyperref}
%\usepackage{todonotes}
\usepackage{caption}
\usepackage{graphicx}
%\usepackage{multirow}
%\usepackage{blindtext}
%\usepackage{abntex2/abntex2-cin-ufpe}
%\usepackage{graphicx}
%\usepackage{array}
%\usepackage[bookmarks=true]{hyperref}
\usepackage{bookmark}
%\usepackage[portuguese]{babel}
%\usepackage[brazil]{babel}
%\usepackage[latin1]{inputenc}
%\usepackage{enumerate}
%\usepackage{cite}
%\usepackage{epsf,epsfig,psfig}
%\usepackage{pagina}
%\usepackage{booktabs}
%\usepackage{indentfirst}
\usepackage{pdflscape}
\usepackage{theorem}
\usepackage{setspace}
\usepackage{boxedminipage}
\usepackage{float}
\usepackage{makeidx}
\usepackage{amsmath}
%\usepackage[]{hyperref}
%\usepackage{todonotes}
%\usepackage{caption}
%\usepackage{multirow}
%\usepackage{blindtext}
%\usepackage{rotating}
%\usepackage[toc,page]{appendix}
\usepackage[alf]{abntex2cite}
\usepackage{comment}
\usepackage{pdfpages}
\usepackage{cmap}
\usepackage[T1]{fontenc}
\usepackage{lastpage}			% Usado pela Ficha catalográfica
\usepackage{tikz}
\usepackage[section]{placeins}
\usepackage{blindtext}
\usepackage{rotating}
\usepackage{adjustbox}
\usepackage[graphicx]{realboxes}
%\usepackage{varioref}
%\usepackage{setspace}
%\usepackage{lipsum}
%\usepackage{scrpage2}
\usepackage{float}
\usepackage{cleveref}
%\usepackage[breaklinks]{hyperref}
\usepackage{prettyref}
%\usepackage[pdftex]{hyperref}
\usepackage[defaultlines=2,all]{nowidow}
%\usepackage[brazilian,hyperpageref]{backref}


\makeindex

\newtheorem{deff}{Definição}[section]
\numberwithin{equation}{chapter}

\theoremstyle{plain}

%\bibliographystyle{abnt-num}

\newcommand*{\captionsource}[3]{%
    \caption{\label{#1}#2.}
    \vspace{-0.3cm}
    \caption*{Fonte: #3.}
}
%\setlength {\marginparwidth }{2cm} 

\begin{document}



\hypersetup{
    citecolor=black,
    filecolor=black,
    linkcolor=black,
    urlcolor=black,
    linktoc=all}
    
% Infos do projeto
%
\newcommand{\theuniversity}{Universidade Federal de Pernambuco}
\newcommand{\thecenter}{Centro de Filosofia e Ciências Humanas}
\newcommand{\thedepartment}{Departamento}
\newcommand{\thecourse}{Curso}
\newcommand{\thetitle}{Título}
%\renewcommand{\anexoname}{ANEXO}
%\renewcommand{\anexosname}{ANEXO}
%\renewcommand{\apendicename}{APPENDIX}
%\renewcommand{\apendicesname}{Appendix}
%\newcommand{\thetype}{Palloma Marciano}
\newcommand{\theproftitle}{}
\newcommand{\thestudent}{Autora}
\newcommand{\theadvisor}{Orientadora}
\newcommand{\thecity}{Recife}


% \thispagestyle{empty}\listoftodos\pagebreak
\thispagestyle{empty}\input{01_Pre_textuais/Capa}
\pagebreak\thispagestyle{empty}\input{./01_Pre_textuais/FolhaDeRosto}
%\pagebreak\thispagestyle{empty}\input{01_Pre_textuais/Projeto}
%\pagebreak\thispagestyle{empty}\input{01_Pre_textuais/Epigrafe}
%\pagebreak\thispagestyle{empty}\input{01_Pre_textuais/Resumo}
%\pagebreak\thispagestyle{empty}\input{01_Pre_textuais/Abstract}

\fancypagestyle{plain}{
\fancyhf{} % clear all header and footer fields
\renewcommand{\headrulewidth}{0pt}
\renewcommand{\footrulewidth}{0pt}}
\pagestyle{plain}

\pagebreak

%\def\listfigurename{LIST OF FIGURES}\listoffigures
%\def\listtablename{LIST OF TABLES}\listoftables
%\input{01_Pre_textuais/ListaDeSiglas.tex}    % não coloquei LISTA DE SIGLAS no meu trabalho, mas fica aqui a chamada comentada.
%\cleardoublepage
%\pdfbookmark{\contentsname}{Contents}
%\tableofcontents
\def\contentsname{SUMÁRIO}\tableofcontents

\fancypagestyle{plain}{
\fancyhf{} % clear all header and footer fields
\fancyhead[R]{\thepage}
\setlength{\voffset}{-1cm}
\setlength{\headsep}{1cm}
\renewcommand{\headrulewidth}{0pt}
\renewcommand{\footrulewidth}{0pt}}

\pagestyle{plain}

\pagebreak
\input{./02_Texto/introduction/introducao.tex}
%\input{./02_Text/introduction/introduction.tex}
%\input{./02_Text/introduction/equacoes_simples.tex}
%\input{./02_Text/introduction/abstract.tex}
\input{./02_Texto/development/theory/object.tex}
\input{./02_Texto/development/theory/theory.tex}
\input{./02_Texto/introduction/objetivos.tex}
%\pagebreak
\input{./02_Texto/development/data_methods/data.tex}

%\input{./02_Text/development/tests/clinical_test.tex}
% inserir demais capítulos aqui
% -----------------------------
% -----------------------------
% -----------------------------
% -----------------------------


\pagebreak
%\input{02_Text/conclusao/conclusao.tex}

%\pagebreak


\def\bibname{REFERÊNCIAS}
\bibliography{Library}
\bibliographystyle{abntex2-alf}
%\addcontentsline{toc}{chapter}{\numberline{}\textbf{REFERÊNCIAS}}


%\input{Anexos/anexo_a} % inclui conte\’udo do arquivo ax01.tex


%\include{Anexos/anexo_1}

%\appendix
%\chapter{Anexo 1}
%\input{anexos}

%\pagebreak
%\addcontentsline{toc}{chapter}{\numberline{}\textbf{ANEXOS}}
%\include{anexo_1}
%\def\appendix{ANEXOS}
%\appendix{./03_Pos_textuais/anexo_1.tex}

%\begin{ANEXO}
%\chapter{ANEXO}
%\end{ANEXO}

%felipe% \printindex    %Removi o índice remissivo para a versão oficial do trabalho.

\end{document}

\end{document}